\origin{ai}

\subsection{Finite forced-window trace-cycle-index kernel}
\label{sec:window6-clock-gauge-invariant-trace-cycle-index-necklace-kernel}

\paragraph{Carrier.}
Fix a positive integer $n$. The finite carrier is a divisor-indexed necklace kernel built from weighted traces. For an indeterminate $z$, set
$$
\begin{aligned}
T(z) &= \begin{pmatrix}1 & z \\ 1 & 0\end{pmatrix},
&
M &= T(1)=\begin{pmatrix}1 & 1 \\ 1 & 0\end{pmatrix}.
\end{aligned}
$$
The matrix $M$ is the unweighted Fibonacci transfer matrix, while $T(z)$ records a weight by attaching one factor of $z$ to the upper-right transition. For each positive integer $g$, the weighted trace is the polynomial $\operatorname{tr}(T(z)^g)$. Its coefficients are the finite numbers
$$
\begin{aligned}
W(g,0)&=1,\\
W(g,s)&=\frac{g}{g-s}\binom{g-s}{s}\quad\text{for }s\geq 1.
\end{aligned}
$$
Thus, for $1\leq g\leq 10$,
$$
\begin{aligned}
\operatorname{tr}(T(z)^g)&=\sum_s W(g,s)z^s.
\end{aligned}
$$
The specialization at $z=1$ returns the trace of the Fibonacci transfer matrix:
$$
\begin{aligned}
\operatorname{tr}(T(1)^g)&=\operatorname{tr}(M^g)=L_g,
\end{aligned}
$$
where, for $g=1,\ldots,10$, the exact list is
$$
\begin{aligned}
(L_1,L_2,L_3,L_4,L_5,L_6,L_7,L_8,L_9,L_{10})
&=(1,3,4,7,11,18,29,47,76,123).
\end{aligned}
$$

\paragraph{Primitive necklace extraction.}
Let $\mu$ denote the integer Mobius function: $\mu(1)=1$; $\mu(r)=0$ when some prime square divides $r$; and, when $r$ is a product of $m$ distinct primes, $\mu(r)=(-1)^m$. For each positive integer $d$, define the primitive weighted necklace polynomial
$$
\begin{aligned}
P_d(z)&=\frac{1}{d}\sum_{r\mid d}\mu(r)\operatorname{tr}(T(z^r)^{d/r}).
\end{aligned}
$$
The trace-cycle-index kernel is then
$$
\begin{aligned}
G_n(u,y)&=\sum_{d\mid n}u^dP_d(y^{n/d}).
\end{aligned}
$$
This is a single finite polynomial in $u$ and $y$. The exponent of $u$ records the cycle divisor $d$, while the exponent of $y$ records the induced weight after the substitution $z=y^{n/d}$. Since all sums are over divisors, the construction is finite and purely symbolic for every fixed $n$.

\paragraph{The case $n=10$.}
For the divisor set $\{1,2,5,10\}$, the primitive polynomials are exactly
$$
\begin{aligned}
P_1(z)&=1,\\
P_2(z)&=z,\\
P_5(z)&=z+z^2,\\
P_{10}(z)&=z+3z^2+5z^3+2z^4.
\end{aligned}
$$
Substituting these four rows into the kernel gives the explicit polynomial
$$
\begin{aligned}
G_{10}(u,y)
&=u+u^2y^5+u^5y^2+u^5y^4+u^{10}y+3u^{10}y^2+5u^{10}y^3+2u^{10}y^4.
\end{aligned}
$$
The joint spectrum is recovered by coefficient extraction:
$$
\begin{aligned}
[u^dy^k]G_{10}(u,y)
&=J_{10}(d,k),
\end{aligned}
$$
with the nonzero values
$$
\begin{aligned}
J_{10}
&=\{(1,0):1,(2,5):1,(5,2):1,(5,4):1,\\
&\qquad (10,1):1,(10,2):3,(10,3):5,(10,4):2\}.
\end{aligned}
$$
This coefficient table is not an additional datum; it is the finite table read directly from the displayed polynomial.

\paragraph{Specializations.}
The same kernel also gives three lower-dimensional summaries. Setting $u=1$ produces the weight polynomial
$$
\begin{aligned}
G_{10}(1,y)&=1+y+4y^2+5y^3+3y^4+y^5,
\end{aligned}
$$
so its coefficient list in increasing powers of $y$ is
$$
\begin{aligned}
(1,1,4,5,3,1).
\end{aligned}
$$
Setting $y=1$ gives the divisor-size spectrum
$$
\begin{aligned}
[u^d]G_{10}(u,1)&=\{1:1,2:1,5:2,10:11\}.
\end{aligned}
$$
Finally, the scalar specialization is
$$
\begin{aligned}
G_{10}(1,1)&=15.
\end{aligned}
$$
Thus the same finite kernel simultaneously generates the joint divisor-weight table, the weight marginal, the divisor-size marginal, and the total orbit count.

\paragraph{Witness extraction.}
The witness is obtained by finite arithmetic over the displayed matrices and divisor sums. One computes $\operatorname{tr}(T(z)^g)$ for the required values of $g$, extracts $P_d$ by the divisor formula, forms $G_{10}(u,y)$, and then reads coefficients and specializations. The bridge back to the unweighted Fibonacci transfer matrix is the identity $T(1)=M$, which gives $\operatorname{tr}(T(1)^g)=\operatorname{tr}(M^g)=L_g$ for the listed range. No limiting argument, asymptotic approximation, or fitted parameter is involved.

\paragraph{Local boundary.}
The construction proves only the finite trace-cycle-index statements displayed above. It does not identify any physical constant, impose a global dynamical law, depend on an external measurement system, compare to metrological data, or assign an interpretation beyond this finite symbolic construction unless an independent certificate supplies those extra assertions.
